\scanpage{23}
Soit maintenant $\beta\in I\setminus K$, prouvons que l'on a $\overline V\cap E_\beta=\varnothing$. En effet, $V$ ne \deleted{contient} \inserted{rencontre} aucun $E_i$ $(i\in I)$ donc $E_\beta\subset X\setminus V=\mathop{\unsure{\deleted{\bigvee}\bigcup|\bigcup}}\limits_{\alpha\in K}\widehat{\deleted{\unsure{I|E}}\Sigma}_\alpha^+$, \deleted{\unsure{donc si}} je dis que si $\beta\in I\setminus K$, on a en fait
\[
E_\beta\subset\bigcup\Sigma_\alpha^+\qquad\text{(ouvert disjoint de }V\text{)}
\]
% slanted
\reversemarginpar\margin{\makebox[\linewidth][r]{\rotatebox{24}{\shortstack[l]{C'est trivial, \deleted{\illegible{1}}\\du fait que les $E_i$ $(i\in\unsure{I|J})$\\sont disjoints}}}}\normalmarginpar
ce qui \unsure{prouvera|prouve} l'assertion -- mieux, je dis que $\exists\alpha\in K$ tel que $E_\beta\subset\Sigma_\alpha^+$. Cela \unsure{revient|mérite} \unsure{au|un}

\uline{Lemme} Soit $K\subset I$, \deleted{tel que les} \deleted{et} des $\Sigma_\alpha^+$ $(\alpha\in K)$ tels que
\[
\alpha\mathrel{\unsure{\ne|\deleted{\ne}{,}}}\beta\in K,\ \alpha\ne\beta\ \Rightarrow\ \Sigma_\alpha^+\cap\Sigma_\beta^+=\varnothing.
\]
Soit $i\in I\setminus K$ tel que $E_i\subset\bigcup\Sigma_\alpha^+$, alors \deleted{$E_i$ est \unsure{contenu|contenue}} \inserted{$\exists\alpha\in K$ tel que} $E_i\subset\Sigma_\alpha^+$.

\unsure{Ce $L^\circ$|Ce $L^o$|Cela} \deleted{signifie que} \unsure{provient|précise|précise le} \unsure{du|des|sens|donc}

\uline{Corollaire} Avec les hyp. préc. sur $(\Sigma_i^\alpha)_{\alpha\in\unsure{K|I}}$,

\noindent\begin{tabular}{@{}|p{.96\linewidth}@{}}
la partition $\Pi_K^{\unsure{\circ|0}}$ est formée des $\Sigma_\alpha^+$ $(\alpha\in K)$ et de $V=X\setminus\mathop{\overline{\bigcup}}\limits_{\alpha\in K}\widehat{\Sigma}_\alpha^+$.
\end{tabular}

Immédiat par récurrence sur $\operatorname{card} K\geq 0$.
% archive sheet number: 22
